\begin{abstract}
\noindent We report on a 67-paper evaluation of AI-agent-assisted replication of
computational science papers drawn primarily from the Argonne National Laboratory
corpus on OSTI.gov, plus select non-OSTI benchmarks in astrophysics,
numerical PDE, generative modeling, microbial genomics, and PDE/PINN methods. The cohort spans 20+ domains
(astrophysics, condensed matter, nuclear physics, computational chemistry,
materials, power systems, computational biology, microbial genomics, climate, fusion, NLP,
control theory, combinatorics, and more). Each replication was executed by an
agentic AI assistant (Ollie / OpenClaw) with tool access and compute on
CherryRd (local), uicgpu (shared 8$\times$A100), Aurora, Polaris, and DGX Spark
nodes. The model backbone is Anthropic Claude Opus via a free Argo proxy at
Argonne, with subagent parallelism enabling 14 new replications in Wave~2
(2026-04-26 to 2026-04-30), 10 biology/BV-BRC replications in Wave~3
(2026-05-05), 5 PDE/PINN replications in Wave~4 (2026-05-07), and 6 additional BV-BRC clinical-genomics replications in Wave~5 (2026-05-10). We scored each attempt on two independent 1--10 axes:
\emph{coverage} (what fraction of the paper's contributions were reproduced)
and \emph{agreement} (how closely quantitative/qualitative results match).
Across the 67-paper cohort, mean coverage is \textbf{7.48/10} and mean
agreement is \textbf{8.02/10}. 35 papers scored $\geq 8$ on both axes; 10
scored $\leq 5$ on at least one axis. The strongest replications occurred
where open-source code plus public data existed (combinatorics, graph
algorithms, topological response, pulsar timing, transit detection,
sequence-level bioinformatics), while the
weakest involved proprietary pipelines (BV-BRC/SEEDtk genome binning), long HPC
campaigns (CROC cosmic reionization), or honest negative results
(Godunov-loss PDE, where the paper's qualitative claim was not reproduced for
the tested MLP architecture). A notable finding from Wave~4 is a systematic
PINN reproducibility gap: 3 of 5 PDE papers using physics-informed neural
networks reproduced qualitatively but showed 10--50$\times$ higher absolute
errors than claimed, attributable to missing author code, hyperparameter
sensitivity, and training details not specified in the papers.
Key infrastructure innovations include:
\texttt{pack\_jobs.py} for HPC batch throughput, subagent-based parallel
replication, BV-BRC API integration for microbial genomics, and a curated
scoring rubric enabling inter-paper comparison.
We identify recurring failure modes, map each underperforming paper to a
concrete upgrade path, and distill 275+ follow-on research questions (5 per
paper) suitable for an AI-assisted research agent. The results support the
thesis that, with appropriate tooling and compute, an AI agent can shoulder
the mechanical burden of replicating a broad swath of published computational
science and surface concrete extensions.
\end{abstract}
